\chapter{Literature Review}
\label{ch:literature}

\section{Z3}
\label{sec:lit-z3}

Z3~\cite{z3_tutorial} is a Satisfiability Modulo Theories (SMT)
solver developed by Microsoft. SMT extends boolean satisfiability
(SAT) by allowing constraints over richer theories such as
bit-vectors, integers and arrays. Internally, Z3 uses CDCL
(Conflict-Driven Clause Learning) to decide satisfiability of
formulas in these theories. DPLL (Davis-Putnam-Logemann-Loveland)
is a classical backtracking algorithm for SAT that combines case
splitting with unit propagation. CDCL extends DPLL by analysing
each conflict during the search and learning a new clause from it,
which prunes the future search by preventing revisits to similar
dead-ends. We used Z3's Python API (Z3py) to encode our
specifications and the Z3 documentation was referred to often
throughout the project.

\section{CEGIS}
\label{sec:lit-cegis}

The Syntax-Guided Synthesis paper~\cite{cegis_paper} and the
Search-Based Program Synthesis primer~\cite{synth_primer} provide
a simple walkthrough of CEGIS. They served as our introduction to
CEGIS and provided the foundations for understanding the
ParserHawk paper. To get more familiar with these concepts, we
built a small CEGIS-based toy DSL which infers a function from a
logical specification. The code for this is present in our GitHub
repository~\cite{parsercompiler26_repo}.

\section{ParserHawk}
\label{sec:lit-parserhawk}

ParserHawk~\cite{parserhawk} is the main reference for our project.
It solves a similar synthesis problem for parsers, but converts P4
code into a TCAM-based implementation. The problem it solves is
different from the problem described before: 
it compiles P4 parsing code to be used in hardware, while we convert
sequential C-like code into a FSM (similar to the output
of ParserHawk) and then back to P4. ParserHawk was very useful in
inspiring ideas which shaped our project, mainly the synthesis
encoding (Section~\ref{sec:synthesis}) and logical specification construction.